Studienabschlussarbeiten sind in aller Regel gebunden einzureichen. Meist
kommt eine Kamm- oder Klebebindung zur Anwendung. Da ein Teil des Innenrandes
von der Bindung und dem gebogenen Papier verdeckt wird, muss der Satzspiegel
ein klein wenig enger gestaltet werden. Somit sind die Proportionen der
aufgeschlagenen Arbeit wieder optisch stimmig. Dieser als Bindekorrektur
bezeichnete Abstand~$K$ bestimmt sich dabei aus der Dicke des Buchblocks~$d$
und der Breite der Bindung~$b$.
%
\begin{equation}
K = \dfrac{1}{2} \cdot \max(d, b)
\end{equation}
%
Die Buchblockdicke lässt sich anhand der Grammatur~$G$ und der Anzahl der
Blätter~$n$ berechnen. Dabei ist die Grammatur in der Einheit $\unit{g / m^2}$
in die Größengleichung einzusetzen. Das Ergebnis der Formel entspricht der
Dicke in Millimetern.
%
\begin{equation}
d = \dfrac{n \cdot G}{1000}
\end{equation}
%
Ein Buchblock aus 100~Seiten (das sind 50~Blätter) hat demnach bei Papier der
Stärke $\unit[80]{g / m^2}$ eine Dicke von \unit[4]{mm}. Eine übliche
Kammbindung hat jedoch eine Breite von \unit[6]{mm}, sodass die Bindekorrektur
auf diesen Wert zu vergrößern ist.

Die Bindekorrektur lässt sich über das Paket \texttt{geometry} einstellen.
Listing~\ref{lst:tpl:bcor} zeigt die entsprechende Konfigurationsoption.

\lstinputlisting[language={[LaTeX]{TeX}},style=block,escapechar=\!,caption={Einstellung der Bindekorrektur},label={lst:tpl:bcor}]{code/tpl_bcor}

Bei Ringbindungen muss keine Bindekorrektur erfolgen.

% DTK 2004/1
